\chapter{Technical Annexes}
\label{cap:anexos}

\section{ODD description of the model}
\label{sec:odd}

The description follows the ODD protocol \cite{grimm2020odd}.

\subsection*{Overview}

\textbf{Purpose.} To study how a constant-product AMM tracks an exogenous reference price after a perturbation, under different arbitrage capacities, depths and liquidity-provision rules.

\textbf{Entities and state variables.} The AMM (reserves $x_t$, $y_t$); the CEX reference ($\pC_t$); four wallets: noise, trend and panic, arbitrageur and liquidity provider, with inventories in $X$ and $Y$. The derived variables are $e_t$, $d_t$ and $z_t$.

\textbf{Scales.} A tick is a discrete unit of time with no calendar equivalence; the horizon is 20 ticks. The initial reserves are $10\,000$ units per asset and the initial price is 1.

\textbf{Process overview.} At each tick: reference shock (only at tick 1), noise order, trend or panic decision, threshold-withdrawal provider rule, arbitrage order, recording and checks.

\subsection*{Design concepts}

\textbf{Emergence.} The trajectory of the tracking error results from the interaction between bounded arbitrage and the order flow. \textbf{Rule-based responses.} The provider withdraws liquidity upon an error greater than $0.02$ (once); the panic agent reacts to an AMM return less than or equal to $-1.5\,\%$ (once per run); the trend agent follows returns of magnitude greater than or equal to $0.1\,\%$. \textbf{Objectives.} None of the agents optimizes; the rules are heuristic. \textbf{Sensing.} The agents observe the AMM price and, the arbitrageur, also the reference. \textbf{Stochasticity.} Only the noise flow is random (activity and side); the AMM mechanics and the rules are deterministic. \textbf{Interaction.} Indirect, through the price and the reserves.

\subsection*{Details}

\textbf{Initialization.} Reserves $10\,000/10\,000$, initial reference price 1; noise and trend wallets with twice the initial reserves, arbitrageur with five times, provider with zero. \textbf{Input.} The reference shock and the noise schedule pre-generated from the seed. \textbf{Submodels.} Equations~\eqref{eq:reservas} to~\eqref{eq:qstar} and the rules of the section on agents in chapter~\ref{cap:metodologia}. Table~\ref{tab:parametros} summarizes the parameters.

\begin{table}[htbp]
\centering
\caption{Parameters of the executed model (version 0.4.0).}
\label{tab:parametros}
\small
\begin{tabular}{@{}lr@{}}
\toprule
Parameter & Value \\
\midrule
Initial reserves $x_0=y_0$ & 10\,000 \\
Initial reference price & 1 \\
Horizon & 20 ticks \\
Shock tick & 1 \\
Band half-width $\gamma$ & 0.003 \\
Ticks inside the band for the sustained return & 2 \\
Arbitrage capacities $c$ & 0; 0.005; 0.03; 0.10 \\
Noise and trend inventory (multiple of the reserves) & 2 \\
Arbitrageur inventory (multiple) & 5 \\
Noise activity probability & 0.7 \\
Noise order size (fraction of the initial reserve) & 0.004 \\
Trend threshold (logarithmic return) & 0.001 \\
Trend order size & 0.006 \\
Panic threshold (logarithmic return) & $-0.015$ \\
Panic sale size & 0.02 \\
Maximum panic sales per run & 1 \\
Threshold-withdrawal trigger (logarithmic error) & 0.02 \\
Fraction of reserves kept & 0.55 \\
\bottomrule
\end{tabular}
\end{table}

\section{Execution contract and reproducibility}
\label{sec:contrato}

Each simulator run is described by the fields of table~\ref{tab:contrato}.

\begin{table}[htbp]
\centering
\caption{Execution contract.}
\label{tab:contrato}
\small
\begin{tabularx}{\textwidth}{@{}l X@{}}
\toprule
Field & Content \\
\midrule
\code{run\_id} & Identifier derived from the parameters. \\
Commit & Code version and state of the working tree. \\
Scenario & Condition, capacity and liquidity rule. \\
Seed & Seed of the noise flow. \\
Parameters & Complete configuration of the model and of the protocol. \\
Clock & Tick and recorded event. \\
Series & Prices, reserves, errors, orders. \\
Metrics & Those of table~\ref{tab:metricas}. \\
Checks & Asset conservation, positivity and finiteness. \\
Environment & Python and platform version. \\
Checksum & SHA-256 of configurations, code and products. \\
\bottomrule
\end{tabularx}
\end{table}

\subsection*{Archived Monte Carlo run record}

\begin{itemize}
  \raggedright
  \item Identifier: \code{\RunID}, analysis version 0.5.0, model version 0.4.0, Python 3.12.3.
  \item Code commit: \ifdefined\blindreviewmode\code{[omitted for anonymous review]}\else\code{06778ccfb4db032d8ae143dcd4396d8786ee7814}\fi, with a clean working tree (\code{git\_working\_tree\_dirty} false).
  \item Run counts: 3000 exploration and 13\,500 held-out simulation runs (historically labeled confirmation).
  \ifdefined\blindreviewmode
  \item Configuration, source, run and product checksums are archived; their values are withheld here because they can identify the public artifacts.
  \else
  \item SHA-256 of the archived confirmation configuration:\newline\code{dcf2ef01fdc7504834362f7cb38b93801b2effa861f6d0e3be8ed74796cef16f}; of the model configuration:\newline\code{9c221ada2b8a2150d238c11d320e34fd3697a6a67d193ab4951b48547b09623f}.
  \item SHA-256 of the source code: \code{confirm.py} \code{9b2c7df09f49bd2f\ldots}, \code{abm.py} \code{f15c666856519b8\ldots}, \code{run.py} \code{1530471be964587e\ldots}, \code{h1c.py} \code{6c26ed25f8796d1c\ldots} (full values in the run summary).
  \item SHA-256 of the products: \code{run\_metrics.csv} \code{2a0fbd17bb3e5eb6\ldots}; \code{inference.csv} \code{d4828a48cd0f8f2a\ldots}; \code{cell\_breakdown.csv} \code{0a87125ca4a09068\ldots}; \code{h3\_by\_condition.csv} \code{28d31f52e5c4be41\ldots};\newline\code{condition\_summary.csv} \code{ee6c02c5ac37544b\ldots};\newline\code{exploration\_fits.csv} \code{1f5f7ffbe00a86c6\ldots}; \code{sample\_size\_guidance.csv} \code{86d4a3a45fedd0a2\ldots}.
  \fi
\end{itemize}

\subsection*{Supplementary engineering sensitivity and full replay}

\ifdefined\blindreviewmode
The original runner sources and archived run configuration were not changed. Five existing pipelines reproduced their recorded output artifacts byte for byte and passed all nine automated implementation and rule checks. Repository-identifying run IDs are withheld in this copy.
\else
The original runner sources and archived run configuration were not changed. Five existing pipelines reproduced their recorded output artifacts byte for byte: deterministic mechanism \code{cex-dex-mvp-f183520b90aa}, sweep \code{cex-dex-sweep-5e33d51c51a3}, stochastic selection \code{cex-dex-abm-a99f9be866b2}, held-out map \code{cex-dex-h1c-43e446f04adf}, and the historically named confirmation run \code{cex-dex-confirm-5b06486d695e}. All nine automated implementation and rule checks passed.
\fi

\ifdefined\blindreviewmode
The supplementary analysis (analysis 0.6.0; model 0.4.0) used 300 fresh seeds (2000--2299) and 40\,500 paths. Its protocol covers three order-flow variants and 15 evaluated cells. Resampling is by seed; cell recovery intervals are Wilson intervals, and paired depth intervals are unadjusted normal approximations. A second run reproduced the numeric CSV and configuration artifacts byte for byte. Repository paths, run identifiers and checksums are withheld in this copy.
\else
The supplementary analysis is \code{cex-dex-engineering-9438e194aa7e} (analysis 0.6.0; model 0.4.0), 300 fresh seeds (2000--2299), 40\,500 paths. Its protocol and configuration are \code{experiments/cex\_dex\_mvp/ENGINEERING\_PROTOCOL.md} and \code{engineering\_config.json}. The independent-unit resampling is by seed; cell recovery intervals are Wilson intervals, and paired depth intervals are unadjusted normal approximations. The run and export checksums are in \code{figures/engineering/manifest.json}. Repeating the analysis in a second output directory reproduced all numeric CSV and config artifacts byte for byte.

\fi
\ifdefined\blindreviewmode
The run products and reproduction commands are withheld in the anonymous copy.
\else
The products of the run are not versioned in the repository. To regenerate them and export the figure data:

\begin{verbatim}
python3 experiments/cex_dex_mvp/confirm.py
python3 experiments/cex_dex_mvp/export_thesis_data.py \
  --run-dir experiments/cex_dex_mvp/output/confirm/<run_id>
python3 experiments/cex_dex_mvp/check_invariants.py
\end{verbatim}
\fi

Repeating the run produced the same identifier and byte-for-byte identical result files, except for the \code{generated\_at\_utc} and \code{git\_working\_tree\_dirty} fields of the summary.

\section{Verification plan}
\label{sec:verificacion}

Table~\ref{tab:verificacion} lists the nine automatic checks of \code{check\_invariants.py}. All passed in the version used for the archived results. The list includes both mechanical implementation checks and checks that expected behaviors follow from the configured rules; passing them is not independent evidence for the reported market effects.

\begin{table}[htbp]
\centering
\caption{Simulator verifications and result.}
\label{tab:verificacion}
\small
\begin{tabularx}{\textwidth}{@{}l X l@{}}
\toprule
Id. & Check & Result \\
\midrule
V01 & The zero-fee sale of $X$ conserves $k$ and keeps reserves positive. & passes \\
V02 & The zero-fee purchase of $X$ conserves $k$ and keeps reserves positive. & passes \\
V03 & Proportional withdrawal conserves the marginal price. & passes \\
V04 & The configured states from E0 to E3 have positive and finite reserves. & passes \\
V05 & The configured arbitrage rule moves the pool price toward the reference for both shock directions. & passes \\
V06 & The fixed withdrawal threshold activates under the executed shocks and applies the configured proportional withdrawal. & passes \\
V07 & The static and threshold-withdrawal comparison shares the exogenous trajectory and the ticks. & passes \\
V08 & The impact of a fixed sale grows with each reduction of liquidity. & passes \\
V09 & Paired conditions share the noise orders and conserve assets (24 trajectories, 504 states). & passes \\
\bottomrule
\end{tabularx}
\end{table}

These checks verify the implementation. Validation against real markets was not performed. There is no automatic check that the simulator reproduces stylized facts, nor was behavior with fees verified, because it is not implemented.

\section{Literature protocol and reference audit}
\label{sec:protocolo-literatura}

The active BibTeX source contains 87 records. The compiled bibliography in this revision contains 56 cited items: 55 scholarly works and one repository record. Metadata checks against publisher records (Crossref) or preprint repositories covered titles, authors, year, publication, volume and pages for the cited works. Scholarly works were prioritized by relevance to the problem: core items support central claims (AMM mechanics, arbitrage, agent-based models and validation), recommended items broaden context, and reserve items support a specific claim. The following cases were identified and are declared to make their status clear:

\begin{itemize}
  \item Campbell et al., Bergault et al., Dean et al. and Cohen et al. are cited as preprints or working papers, with no peer review located.
  \item In Dean et al., the Crossref record lists Dean as first author, and another bibliographic database lists a different name; Crossref is used and it is recommended to confirm it on the publisher's page of the working paper.
  \item Windrum et al. (2007) was located only by URL and has no DOI verified in the record consulted.
  \item Schär (2021): the pages were inferred from the record and not confirmed on the journal's site. Angeris et al. (2020): volume and issue were not confirmed.
\end{itemize}

The detailed list, with the evidence per entry, is kept \ifdefined\blindreviewmode in the study materials.\else in the repository (\code{thesis/literature-audit/}).\fi\ No earlier ABM of an AMM with liquidity withdrawal dependent on the tracking error was found; this claim is limited to the databases consulted.

\section{Sensitivity protocol}
\label{sec:sensibilidad}

Further sensitivity work should vary the candidate factors in table~\ref{tab:factores}. One-factor-at-a-time changes can screen individual contributions; factorial or screening designs \cite{morris1991factorial,saltelli2010variance} are preferable when interactions among several factors are the target. Before running a new experiment, specify its primary estimand, contrast, seed-level resampling unit, practical precision target and stopping rule, then record the sample-size rationale with the protocol and outputs. The archived \code{figures/data/sample\_size\_guidance.csv} is tied to an earlier design and is not a generic sample-size guide.

\ifdefined\blindreviewmode
The H2 mechanism analysis in section~\ref{cap:resultados} replays 9,000 paths using existing random seeds. It includes matched panic/no-panic paths and a shallow-depth capacity-reference control; 7,200 reserve-linked-cap paths were checked against archived metrics. The replay is exploratory, does not provide an independent sample or external validation, and does not identify within-tick causality. Repository paths and run identifiers are withheld in this anonymous copy.
\else
The H2 mechanism analysis in section~\ref{cap:resultados} can be regenerated with \code{python3} \url{experiments/cex_dex_mvp/analyze_h2_revision.py}. Replay \code{\HtwoReviewRunID} uses 9,000 paths with existing random seeds, including matched panic/no-panic paths and a shallow-depth capacity-reference control. The script writes cell-level CSV files and a manifest with input and output checksums; 7,200 reserve-linked-cap paths are checked against archived metrics. The analysis is exploratory, not an independent sample or external validation, and end-of-tick episodes do not identify within-tick causality.
\fi
